\appendix
\chapter{精选习题与答案骨架}

\section{固体物理}
\begin{enumerate}
\item \textbf{倒格体积。}证明原胞体积满足
$\Omega_c\Omega_c^*=(2\pi)^3$。

\textbf{答案骨架：}用 $\vb b_i$ 的交叉乘积定义计算
$\vb b_1\cdot(\vb b_2\times\vb b_3)$；或把基矢组成矩阵
$A$，由 $B^TA=2\pi I$ 得
$\det B=(2\pi)^3/\det A$。

\item \textbf{fcc 消光。}推导同种原子的 fcc 结构因子选择定则。

\textbf{答案骨架：}取常规胞四个格点
$0,(0,\frac12,\frac12),(\frac12,0,\frac12),
(\frac12,\frac12,0)$，求
$S=f[1+(-1)^{k+l}+(-1)^{h+l}+(-1)^{h+k}]$；全奇或全偶时为 $4f$，其他为零。

\item \textbf{一维态密度。}由
$E=-2t\cos ka$ 求每长度态密度。

\textbf{答案骨架：}$g(E)=2(L/2\pi)|\dd k/\dd E|$，因同一能量有 $\pm k$；
\[
\frac{g(E)}L=\frac1{\pi a\sqrt{4t^2-E^2}},
\]
带边平方根发散。

\item \textbf{Debye 标度。}若 $d$ 维玻色模色散
$\omega=Aq^s$，求低温热容幂律。

\textbf{答案骨架：}$N(q)\propto q^d$，故
$g(\omega)\propto\omega^{d/s-1}$；
$U\propto T^{1+d/s}$，所以
$C\propto T^{d/s}$。

\item \textbf{两带 Hall。}说明电子、空穴完全补偿时仍可有大磁阻。

\textbf{答案骨架：}写两带电导张量并求逆；$n=p$ 时 Hall 项可相消，而纵向轨道弯曲不能相消，理想模型磁阻约随
$\mu_e\mu_hB^2$ 增长。
\end{enumerate}

\section{量子与统计}
\begin{enumerate}
\item \textbf{Kramers 正交。}严格证明
$\langle\psi|\Theta\psi\rangle=0$。

\textbf{答案骨架：}利用反幺正内积性质
$\langle\Theta\phi|\Theta\psi\rangle=\langle\psi|\phi\rangle$，
令 $\phi=\Theta\psi$ 并用 $\Theta^2=-1$。

\item \textbf{黄金规则量纲。}检查
$2\pi|V_{fi}|^2\rho(E)/\hbar$ 的量纲。

\textbf{答案骨架：}$|V|^2$ 为能量平方，$\rho$ 为每能量，除以
$\hbar$ 后为每时间。

\item \textbf{涨落热容。}从 $Z$ 二阶导数推导
$\langle(\Delta E)^2\rangle=\kb T^2C_V$。

\textbf{答案骨架：}$\partial_\beta^2\ln Z=\langle E^2\rangle-\langle E\rangle^2$，
再用 $\partial_\beta=-\kb T^2\partial_T$。

\item \textbf{Landau 外场。}在 $T=T_c$ 求 $m(h)$。

\textbf{答案骨架：}极值条件 $bm^3=h$，所以
$m\propto h^{1/3}$，平均场临界指数 $\delta=3$。

\item \textbf{二次量子化。}证明粒子数算符
$N=\sum_\alpha a_\alpha^\dagger a_\alpha$ 满足
$[N,a_\beta^\dagger]=a_\beta^\dagger$。

\textbf{答案骨架：}使用对易/反对易代数逐项计算；结果表明产生算符把粒子数提高一。
\end{enumerate}

\section{量子材料}
\begin{enumerate}
\item \textbf{Hubbard 二聚体。}在 $U\gg t$ 下比较单态和三重态二阶能量。

\textbf{答案骨架：}三重态受 Pauli 阻止；单态耦合两个双占据中间态。二阶能量差
$E_T-E_S=4t^2/U$。

\item \textbf{Dirac Berry 曲率。}对
$H=v(k_x\sigma_x+k_y\sigma_y)+m\sigma_z$ 判断曲率在何处最大、质量变号有何影响。

\textbf{答案骨架：}两带公式给
$\Omega_\pm\propto\mp mv^2/(v^2k^2+m^2)^{3/2}$，集中于
$k=0$；$m$ 变号使曲率反号并改变正则化后的拓扑贡献。

\item \textbf{SSH 端态。}开边界、零能时求端态振幅递推。

\textbf{答案骨架：}取只占一个子格的解，递推
$\psi_{A,n+1}=-(t_1/t_2)\psi_{A,n}$；当
$|t_1/t_2|<1$ 时从左端指数局域。

\item \textbf{BCS 态密度。}由
$E=\sqrt{\xi^2+\Delta^2}$ 推超导态密度。

\textbf{答案骨架：}守恒态数
$N_s(E)\dd E=N_0\dd\xi$，故
\[
N_s(E)=N_0\frac{|E|}{\sqrt{E^2-\Delta^2}}
\Theta(|E|-\Delta),
\]
带边出现相干峰。

\item \textbf{屏蔽。}由
$V(q)\propto1/(q^2+q_{\rm TF}^2)$ 判断实空间形式。

\textbf{答案骨架：}三维 Fourier 变换得到 Yukawa 势
$V(r)\propto e^{-q_{\rm TF}r}/r$，屏蔽长度为
$q_{\rm TF}^{-1}$。
\end{enumerate}

\section{实验读图}
\begin{enumerate}
\item 给定 $C/T$ 对 $T^2$ 直线，说明如何取得 $\gamma,\Theta_D$。
\item 给定不同光子能量的 ARPES 切图，怎样识别 $k_z$ 色散？
\item 给定 $\rho_{xy}(B)$ 非线性曲线，为什么不能直接用单带 $1/ne$？
\item 给定超导 STS 零偏压峰，列出至少三种替代解释和判别实验。
\item 给定量子振荡不同温度振幅，如何拟合回旋质量？
\end{enumerate}
\textbf{答案共同骨架：}先写直接观测量与模型表达式；再列必要拟合参数和极限；最后至少给一个系统误差与一个替代机制。

\chapter{符号与易混概念索引}
\begin{longtable}{p{0.19\textwidth}p{0.33\textwidth}p{0.38\textwidth}}
\toprule
对象&定义&不要混淆\\\midrule
$\vb k$&Bloch 准动量标签&不等于电子机械动量\\
$\vb G$&倒格矢&不等于任意散射动量\\
$m^*$&能带曲率定义的响应参数&不等于静止质量变化\\
$A(\vb k,\omega)$&单粒子谱函数&不等于 ARPES 原始强度\\
$g(E)$&态密度&不等于占据数 $g(E)f(E)$\\
$\tau_{\rm tr}$&输运寿命&不一定等于量子寿命或谱峰寿命\\
$\Delta$&超导序参量/谱隙能标&复杂体系中两者不必简单相等\\
$\xi$&相干或关联长度&不同章节定义和物理对象不同\\
$\mathcal A$&Berry 连接&依赖规范，曲率或闭合相位才直接物理\\
$C$&Chern 数&需对完整隔离带或占据子空间定义\\
$U$&Hubbard 局域排斥&不等于所有 Coulomb 相互作用\\
$J$&交换常数&符号依 Hamiltonian 约定\\
$\rho$&可指密度、电阻率或密度矩阵&回答时必须先声明\\
$\chi$&广义响应函数&需注明响应量、外场、频率和波矢\\
$\mu$&化学势或磁矩符号&由上下文区分\\\bottomrule
\end{longtable}

\chapter{最后一周复习清单}
\begin{enumerate}
\item 能闭卷推导 Bloch 定理、区边界开隙和紧束缚色散。
\item 能从模式计数推导 Debye $T^3$，从费米面窗口解释电子 $T$。
\item 能写 Boltzmann 方程并说明电导中的速度、寿命和费米面权重。
\item 能从正则系综推出自由能、平均能量、涨落与响应。
\item 能区分能带绝缘、Mott 绝缘、Anderson 局域化。
\item 能由 Hubbard 二阶过程推出 $J=4t^2/U$。
\item 能解释 Berry 连接、曲率、Chern 数和体边对应。
\item 能用 SSH、Dirac 和石墨烯三个最小模型讲拓扑。
\item 能写 BCS Nambu 矩阵、Bogoliubov 谱和能隙方程骨架。
\item 能说明 DFT 自洽循环及 Kohn--Sham 能带的解释边界。
\item 能说清 ARPES、STM、衍射、输运、比热分别测什么。
\item 能对每张数据图提出至少一个替代解释。
\item 能把本科科研表述为“问题--方法--我的贡献--验证--局限”。
\item 对 30 道示范题各完成一次 30 秒和 2 分钟录音。
\item 睡眠、材料和设备提前准备；面试当天不再学习全新专题。
\end{enumerate}

